% Conclusions: restate the contribution and the forward path. No table.

An autonomous Claude Code Opus 4.7 (1M context) Max run built, executed, and
documented a standard library oncology automation platform under strict
assurance, and the evidence supports the thesis directly. The pipeline produced a
complete, well formed deliverable in well under a millisecond, while the VVUQ
gate accepted only one of six candidates and blocked the other five, one of them
escalated to a human. Read against the three questions posed in the introduction,
the platform answered all of them in the right order: verification confirmed that
the code was built correctly, the gate is where validation asks whether the right
deliverable was built, and uncertainty quantification bounds how confident the
result is across repeated runs. The generation answered none of these on its own.
The assurance process, not the generation, is therefore the substantial and
decision bearing part of the workflow, and holding VVUQ to a higher standard than
code generation is precisely what makes physical AI oncology trial deliverables
faster, less expensive, and more rigorous than conventional verification.

The forward path follows from the same logic. The gate that cleared software
deliverables in Stage 1 is wired as an interlock into the Stage 2 2030 pancreatic
cancer analog, so the assurance discipline carries from code into physical safety
without a change of principle. With an independent validation reference supplied
and the real input corpus wired in, the same gate extends to oncology clinical
trials broadly, letting the field check these assurance steps off once rather
than solving them repeatedly, with patient benefit preserved throughout. That is
the contribution of this work: not a faster generator, but an assurance layer
held deliberately above the generator, automated end to end, and ready to scale
~\cite{repo-cancer-automated}.
